% mage-auto-heading: uber
% mage-auto-heading: verso-evidence
% mage-auto-heading: recto-connection-to-mage
% Authoritative LaTeX manuscript; Markdown is a generated compatibility projection.

\subsection{Uber}\label{heading:uber}

\emph{Ride-hailing and logistics platform · a durable environment around a replaceable reasoner (economics motive)}

\subsubsection{Verso --- Evidence}\label{heading:verso-evidence}

\textbf{A shared concern, two motives.} GitLab argues that project knowledge and controls should remain useful when the foundation model or agent changes.\citep{staples2026abundant} Uber reports how it measures and improves the environment around its agents. The two accounts approach a shared engineering concern from different directions: GitLab emphasizes trust in generated work, while Uber emphasizes operating cost. Uber supplies the fuller operational account, which this card examines.

\textbf{What the public record shows.} Uber describes AI tools operating throughout the software lifecycle: more than 70 percent of pull requests are attributed to local or cloud agents, engineers have created more than 3,600 agent skills, and those skills execute more than 30,000 times per day.\citep{medisetty2026factory} Managed agents perform code review, CI repair, end-to-end changes, alert triage, debugging, and maintenance. The interesting structure is not the volume of agent use but that Uber engineers the environment around the agents: real workloads become benchmarks; models are selected against cost, quality, and reliability; context and tool descriptions are supplied on demand to reduce unnecessary reasoning; recurring workflows are recorded as reusable skills; and agents are evaluated in units such as cost per merged pull request, review, or alert.

\textbf{Boundary of the evidence.} Uber reports operational measurements of context, skills, model selection, and cost. These are company reports, not a controlled comparison. They do not establish that executable engineering models guide implementation. GitLab's cited account presents an architectural argument rather than measured operating results.

\textbf{Portable lesson.} Invest in the project knowledge, tools, and checks that later agents can reuse, even when the underlying foundation model changes.

\subsubsection{Recto --- Connection to MAGE}\label{heading:recto-connection-to-mage}

\begin{figure}[htbp]
\centering
\includegraphics[width=\linewidth]{book/latex/assets/field-guide-uber-ea237e11ec.pdf}
\caption{\emph{Uber's connection to MAGE.} Context selection, a knowledge graph, and reusable skills supply agents with project knowledge. The cited account does not establish an executable model of software behavior.}
\label{field-guide-uber}
\end{figure}

\textbf{Connection to MAGE.} Uber illustrates why engineers invest in the environment around an agent. Its tools supply relevant context on demand, its knowledge graph records organizational relationships, and its reusable skills preserve workflows --- engineering capital that can remain useful when the underlying foundation model changes.

\textbf{What the evidence establishes.} Uber describes context selection, a queryable knowledge graph, and reusable skills. These support the connection to MAGE's Modeling practices. The cited account does not establish an executable model of software behavior or a check that rejects disagreement between such a model and the implementation. It also does not measure an increase in reasoning horizon.
